\section{A normal subtraction and all generalized Stieltjes jets}
\label{sec:stieltjes}

\subsection{The correct escape amplitude}
Write
\[
 \Z_{r,N}(s,a)=\sum_{0\le n_1<\cdots<n_r\le N}
       \Delta(\bm n)^2\prod_iw_{n_i}(s,a),\qquad \Z_{r,-1}=0,
\]
and define the largest-index block
\begin{equation}\label{eq:block}
 B_N(s,a)=\Z_{r,N}(s,a)-\Z_{r,N-1}(s,a).
\end{equation}
It vanishes for $N<r-1$. For $N\ge1$, finite Cauchy--Binet on the
remaining particles gives
\begin{equation}\label{eq:block-det}
 B_N(s,a)=w_N(s,a)
 \det\left(\sum_{n=0}^{N-1}n^{i+j}(N-n)^2w_n(s,a)\right)_{0\le i,j<r-1}.
\end{equation}
The amplitude that must be subtracted is
\begin{equation}\label{eq:amplitude}
 A_r(s,a)=c_d(s)V_r(s,a),\qquad c_d(s)=\frac{(s)_d}{d!},
\end{equation}
not just $V_r(s,a)$ away from $s=1$.

\begin{lemma}[Uniform block remainder]\label{lem:block}
Locally uniformly for $\Rea s>0$, $\Rea a>0$,
\begin{equation}\label{eq:block-asymptotic}
 B_N(s,a)=A_r(s,a)(N+a)^{-s}
       +O\bigl((N+1)^{-\Rea s-1}\bigr).
\end{equation}
In particular the difference and all its fixed parameter derivatives
are locally normally summable in that domain.
\end{lemma}
\begin{proof}
The $(r-1)$-particle ensemble with the same weight has total mass $V_r(s,a)$
for $\Rea s>-1$: its largest required moment has degree $d-2$, and
\eqref{eq:moment} without the zeta term extends to that domain by the
identity theorem. For the present proof it suffices to work on a compact
set with $\Rea s\ge\sigma>0$.

The absolute ensemble with one insertion of $\sum_i n_i$ is uniformly
summable. Indeed, bounding its Vandermonde as in Lemma~\ref{lem:decay},
the largest exponent after the insertion is at most $-\sigma-1$.
For $0\le n_i<N$,
\[
 0\le1-\prod_{i=1}^{r-1}(1-n_i/N)^2
       \le\frac2N\sum_{i=1}^{r-1}n_i.
\]
The sum of the absolute ensemble weights with some $n_i\ge N$ is also
at most $N^{-1}$ times the same first-moment majorant. Therefore
\begin{equation}\label{eq:lower-normalization}
 N^{-d}\frac{B_N(s,a)}{w_N(s,a)}=V_r(s,a)+O(N^{-1}),
\end{equation}
where the quotient notation means the lower-particle polynomial sum,
so no division by a zero weight is needed. This proof uses absolute
majorants and is valid for complex $s,a$ as well.

Combine \eqref{eq:lower-normalization} with
\eqref{eq:weight-asymp} and
$N^d(N+a)^{-d}=1+O(N^{-1})$ to obtain
\eqref{eq:block-asymptotic}. Every bracket is holomorphic. Cauchy
estimates on a slightly larger compact neighborhood give summable
majorants for each fixed derivative; equivalently one can retain the
resulting powers of $\log(N+1)$ in the estimate.
\end{proof}

\begin{theorem}[Ordinary pole-cancelled series]\label{thm:subtraction}
The bracketed series
\begin{equation}\label{eq:T}
 T_r(s,a)=\sum_{N\ge0}\left\{B_N(s,a)-A_r(s,a)(N+a)^{-s}\right\}
\end{equation}
is locally normally convergent for $\Rea s>0$, $\Rea a>0$, and
\begin{equation}\label{eq:Tclosed}
 \boxed{\displaystyle
 T_r(s,a)=U_r(s,a)+(1-c_d(s))V_r(s,a)\zeta(s,a).}
\end{equation}
The apparent pole at $s=1$ is removable. The right side gives an entire
continuation in $s$; the assertion that the series itself is ordinary
and normally convergent is limited to the stated half-plane.
\end{theorem}
\begin{proof}
Normal convergence is Lemma~\ref{lem:block}. For $\Rea s>1$, both
unsubtracted series converge, so Theorem~\ref{thm:single} gives
\eqref{eq:Tclosed} by subtraction. The identity theorem then gives the
whole half-plane $\Rea s>0$. Since $c_d(1)=1$, the unique pole of
Hurwitz zeta cancels. The functions $U_r,V_r,c_d$ are entire in $s$.
\end{proof}

\subsection{An explicit triangular system for every Stieltjes constant}
Let
\begin{equation}\label{eq:jetdefs}
 U_k=\partial_s^kU_r(1,a),\quad
 A_k=\partial_s^kA_r(1,a),\quad
 b_k=\left.\partial_s^k\{(1-c_d(s))V_r(s,a)\}\right|_{s=1}.
\end{equation}
Thus $b_0=0$. Also write $T_k=\partial_s^kT_r(1,a)$.

\begin{theorem}[All-order Stieltjes identities]\label{thm:jets}
For every integer $k\ge0$,
\begin{align}
 T_k
 ={}&\sum_{N\ge0}\left\{\partial_s^kB_N(1,a)
   -\frac1{N+a}\sum_{j=0}^k\binom kj A_j
                   [-\log(N+a)]^{k-j}\right\}\label{eq:Tjets-series}\\
 ={}&U_k+\frac{b_{k+1}}{k+1}
   +\sum_{j=1}^k\binom kj b_j(-1)^{k-j}\gamma_{k-j}(a).
 \label{eq:Tjets-closed}
\end{align}
The series is an ordinary locally normally convergent bracketed series.
For $\Rea a>0$, it can be solved recursively for every $\gamma_m(a)$ because
\begin{equation}\label{eq:b1}
 b_1=-H_dV_r(1,a)\ne0.
\end{equation}
\end{theorem}
\begin{proof}
Differentiate \eqref{eq:T} under its normally convergent sum; Leibniz's
rule yields \eqref{eq:Tjets-series}. Set $s=1+u$ in
\eqref{eq:Tclosed}. Multiply
\[
 \sum_{j\ge1}\frac{b_j}{j!}u^j
 \quad\hbox{by}\quad
 \frac1u+\sum_{m\ge0}\frac{(-1)^m\gamma_m(a)}{m!}u^m.
\]
The pole contributes $b_{k+1}/(k+1)$ to the $k$th derivative. The
remaining product contributes the finite convolution in
\eqref{eq:Tjets-closed}. Since
$c_d'(1)=\sum_{j=1}^d1/j=H_d$, equation~\eqref{eq:b1} follows from
Theorem~\ref{thm:single}.
\end{proof}

In explicit recursive form, putting $k=m+1$ gives
\begin{align}
 \gamma_m(a)=\frac{(-1)^m}{(m+1)b_1}
 \left\{T_{m+1}-U_{m+1}-\frac{b_{m+2}}{m+2}\right.\qquad\notag\\
 \left.{}-\sum_{j=2}^{m+1}\binom{m+1}{j}
       b_j(-1)^{m+1-j}\gamma_{m+1-j}(a)\right\}.
 \label{eq:gamma-invert}
\end{align}
The empty sum at $m=0$ is zero. Every coefficient in this triangular
system is a finite combination of powers and logarithms at $a,a+1,\ldots,a+d-1$
and of ordinary harmonic numbers of length $d$.

The amplitude derivatives have a compact all-order description:
\begin{equation}\label{eq:cd-bell}
 \log c_d(1+u)=\sum_{q\ge1}\frac{(-1)^{q-1}}qH_d^{(q)}u^q,
 \qquad |u|<1.
\end{equation}
Consequently
\[
 c_d^{(k)}(1)=\Bell_k\bigl(H_d,-H_d^{(2)},2!H_d^{(3)},\ldots,
                  (-1)^{k-1}(k-1)!H_d^{(k)}\bigr).
\]
Leibniz's rule with the finite determinants $U_r,V_r$ now generates
all the rows without a new infinite summation problem.

\subsection{Concrete low-order identities}
For $r=2$, $a=1$, put $L=\log2$. The finite coefficient functions are
\[
 U(s)=-\frac12+\frac{2^{-s}}4,\qquad
 V(s)=\frac{1-2^{-s}}2,\qquad c_2(s)=\frac{s(s+1)}2.
\]
Equations~\eqref{eq:Tjets-series}--\eqref{eq:Tjets-closed} give
\begin{align}
 T_0&=-\frac34,\label{eq:T0example}\\
 T_1&=-\frac18-\frac12L-\frac38\gamma,\label{eq:T1example}\\
 T_2&=\frac12L^2-\frac14L-\frac{1+3L}{4}\gamma
          +\frac34\gamma_1.\label{eq:T2example}
\end{align}
For example, $A_0=1/4$ and $A_1=3/8+L/4$, so the second line means
precisely
\begin{equation}\label{eq:T1seriesexample}
 \sum_{N\ge0}\left\{B_N'(1,1)
 -\frac{3/8+L/4-\tfrac14\log(N+1)}{N+1}\right\}
       =-\frac18-\frac12L-\frac38\gamma.
\end{equation}
Here $B_N'$ is computed by differentiating the finite two-particle block,
using \eqref{eq:order-weight}. The cancellation belongs to each displayed
bracket; the two component series must not be split at $s=1$.

\subsection{Hard cutoffs are not spectral finite parts}
\begin{corollary}[The cutoff constant and the amplitude derivative]
\label{cor:cutoff}
For $a>0$, put $U_0=U_r(1,a)$ and $V_0=V_r(1,a)$. Then
\begin{equation}\label{eq:hardcutoff}
 \lim_{N\to\infty}\left\{\Z_{r,N}(1,a)-V_0\log(N+a)\right\}
       =U_0+V_0\{\gamma_0(a)-H_d\}.
\end{equation}
The spectral Laurent constant is
\begin{equation}\label{eq:spectralFP}
 \FP_{s=1}\Z_r(s,a)=U_0+\partial_sV_r(1,a)+V_0\gamma_0(a).
\end{equation}
The spectral constant minus the hard-cutoff constant equals
\begin{equation}\label{eq:FPdifference}
 \partial_sV_r(1,a)+H_dV_0=\partial_sA_r(1,a).
\end{equation}
\end{corollary}
\begin{proof}
The finite sum of the brackets in \eqref{eq:T} at $s=1$ is
\[
 \Z_{r,N}(1,a)-V_0\sum_{n=0}^N(n+a)^{-1}.
\]
It tends to $T_0=U_0-H_dV_0$. The digamma shift identity gives
\[
 \sum_{n=0}^N(n+a)^{-1}=\psi(N+a+1)-\psi(a),
\]
and $\psi(N+a+1)-\log(N+a)\to0$. This proves
\eqref{eq:hardcutoff}. Expanding \eqref{eq:single} at $s=1$ gives
\eqref{eq:spectralFP}, and subtraction proves \eqref{eq:FPdifference}.
\end{proof}

For $r=2,a=1$, the constants are respectively
\begin{equation}\label{eq:cutoff-example}
 \frac\gamma4-\frac34,
 \qquad
 \frac\gamma4-\frac38+\frac{\log2}{4}.
\end{equation}
Their difference is $3/8+\log2/4$. The earlier determinantal report
already warns that residue-amplitude derivatives affect this comparison
\cite{repo:dhr}; the present calculation implements that warning for the
new finite-difference spectral deformation. It is not a claim of an error
in that source.

\subsection{Stieltjes primitives with fixed normalization}
For any $m\ge0$ and $\Rea a,\Rea a_0>0$,
\begin{equation}\label{eq:stieltjes-primitive}
 \int_{a_0}^a\gamma_m(x)\dd x
 =\frac{(-1)^{m+1}}{m+1}
    \left\{\zeta^{(m+1)}(0,a)-\zeta^{(m+1)}(0,a_0)\right\}.
\end{equation}
Indeed, comparing coefficients at $s=0$ in
$\partial_a\zeta(s,a)=-s\zeta(s+1,a)$ proves the derivative of
\eqref{eq:stieltjes-primitive}; the anchor fixes its constant. This
classical Hurwitz-jet primitive is consistent with the canonical
manuscript's conventions \cite{repo:integration,dlmf:hurwitz}.
For $m=0$, Lerch's identity reduces it to
\[
 \int_{a_0}^a\gamma_0(x)\dd x=-\log\Gamma(a)+\log\Gamma(a_0).
\]
Substituting \eqref{eq:gamma-invert} supplies ordinary determinant-series
representatives for the integrands of these primitives. Local normal
convergence permits termwise integration along every compact path in
the right half-plane. Thus the Stieltjes series, Gamma primitive, and
spectral derivative descriptions have one consistent normalization.
